% Literature Review: Causal Loop Diagram (CLD) Generation with Large Language Models
% Background section on LLM-based approaches to automated CLD construction

\section{LLM-Based Causal Loop Diagram Generation}
\label{sec:lit_cld_generation}

Notably, generating causal loop diagrams directly is relatively under-researched in the literature found in this review, with most approaches focusing on giving the LLM text, or a causal graph, from which the causal structure is to be extracted or reasoned over, with hallucination mitigation methods largely coming down to different forms of retrieval-augmented generation. To understand the capabilities and limitations of LLMs for CLD-related tasks, we first survey benchmark studies that evaluate different forms of causal reasoning, before examining specific CLD generation and extraction approaches.

\subsection{From LLM-Powered Causal Pattern Matching to Causal Loop Diagrams}

Benchmark studies reveal highly varying performance across causal task types using LLMs, though each evaluates a distinct form of causal reasoning:

\textbf{COPA} (Choice of Plausible Alternatives) tests commonsense causal judgement: given a premise (e.g., ``The man broke his toe'') and two alternatives, models select which is more plausibly the cause or effect. This multiple-choice format with everyday scenarios yields near-perfect GPT-4 accuracy but requires no graph construction or variable extraction---it tests recognition of familiar causal patterns rather than systematic causal inference.

\textbf{CLadder}~\citep{jin2023cladder} tests causal query answering over provided graphs: models receive a causal graph with binary variables and must answer queries at three levels of Pearl's ladder---associational (``What is P(Y|X)?''), interventional (``What happens if we set X=1?''), and retrospective what-if (``Would Y have occurred if X had been different?''). GPT-4 achieves only $\sim$62\% accuracy, with performance degrading sharply on interventional and retrospective what-if queries. \textit{Relation to CLD generation}: CLadder evaluates reasoning \textit{with} a given causal structure, whereas CLD generation requires synthesizing a plausible causal structure from a task prompt or domain specification without being given the graph. However, polarity errors in CLD generation may reflect similar failures in interventional reasoning.

The \textbf{CaLM benchmark}~\citep{chen2024calm} provides the most comprehensive evaluation, testing 28 models across 92 causal targets organised into four modules. Tasks range from causal attribution (identifying causes of events) to effect prediction and what-if reasoning. The key finding---``as complexity of causal reasoning increases, accuracy progressively deteriorates, eventually falling almost to zero''---suggests that CLD generation involving multi-step causal chains will face compounding errors. This is especially relevant for CLDs, where feedback loops and longer cause--effect chains mean that errors in individual edges can propagate through downstream reasoning about the system. It also motivates edge-level error correction to prevent compounding failures, for example via LLM-as-a-judge and LLM-as-a-corrector pipelines (if these components are empirically shown to work in this setting).

\textbf{CausalBench} evaluates three subtasks: correlation identification, causal skeleton identification (undirected graph structure), and causality identification (directed edges). LLMs significantly lag behind traditional algorithms on networks exceeding 50 nodes and struggle with collider structures while excelling at chains.

\citet{darvariu2024largelanguagemodelseffective} frame LLMs as low-cost sources of background knowledge for causal graph discovery: by eliciting structural priors from an LLM, the hypothesis space over candidate graphs can be reduced before running downstream discovery algorithms. They propose metrics to assess LLM causal judgements independently of any specific discovery method and show that LLM-derived priors are most useful for edge directionality (e.g., 84\% correctly oriented edges for LLaMA2-70B when combined with observational data).

\textbf{Corr2Cause}~\citep{jin2024largelanguagemodelsinfer} tests inferring causal graph structure from correlation statements: given statements like ``A is correlated with B'' and ``B is correlated with C,'' models must determine whether ``A causes B'' or identify the correct DAG structure. Testing 17 models on 200K+ samples, all achieved near-random performance ($\sim$29\% F1). This benchmark isolates pure causal structure inference without natural language complexity. The near-random performance suggests low success for text-to-CLD systems in settings where causal structure must be inferred from limited signals. In our CLD generation task, we similarly prompt the model to answer ``Does A cause B?'' (surrounded by prompt scaffolding) at the edge level; based on Corr2Cause, this edge-level decision may be challenging, although this remains to be seen since Corr2Cause conditions on correlation statements rather than the domain-grounded prompts used in this thesis.

However, evidence closer related to our use case shows that, for causal graph discovery from domain knowledge, LLMs show complementary strengths and weaknesses. \citet{jiralerspong2024efficientcausalgraphdiscovery} demonstrate that querying LLMs about pairwise causal relationships using breadth-first search can achieve state-of-the-art results on real-world causal graphs (typically DAG-oriented discovery). This is adjacent to CLD generation in that it elicits edge-level causal hypotheses from domain knowledge, but it does not, by itself, target signed feedback-loop structures as CLDs do.

In the text-to-CLD causal extraction task, an example is the System Dynamics Bot~\citep{hosseinichimeh2024textmapdynamicsbot}, which extracts causal variables and directed causal links from natural language text to construct feedback diagrams and outputs structured CLDs. Evaluated against human-coded ground truth from system dynamics literature and participant responses, the bot recovers only 59\% of causal links and 66\% of feedback loops using GPT-4.

\citet{liu2024leveraging} use curated prompting techniques extracting CLDs from dynamic hypothesis texts, comparing four prompt designs: zero-shot, few-shot, guided instructions, and a two-stage prompt that first extracts variables then maps signed causal links. On 44 textbook DH--CLD pairs (expert CLDs hand-encoded in DOT), curated prompting---especially the two-stage setup, which the generator algorithm employed in this thesis also uses---produces CLDs closest to expert diagrams for simple structures, but errors persist (missed loops, wrong polarities, and weak handling of exogenous/implicit relationships) as complexity grows.

Complementing these point evaluations, the SD-AI benchmark~\citep{schoenberg2025aibuildsdmodels} proposes a standardised test suite and public leaderboard for CLD \emph{causal translation} (i.e., text-to-CLD extraction) and conformance, showing substantial variation across LLMs and motivating continued benchmarking of instruction-following and technical correctness in AI-assisted system dynamics tooling.

An approach closer to home exploring CLD generation comes from \textbf{Theoraizer}~\citep{Waaijers2024theoraizer}, which is closest to the approach we explore in this thesis. Theoraizer operationalises CLD construction as a stepwise, edge-level elicitation pipeline in which an LLM functions as a ``digital extension'' of an expert group rather than a replacement. Users provide (or optionally generate) a variable list, and the system then queries the LLM over all ordered variable pairs to (i) classify whether a causal relation exists, (ii) infer directionality, and (iii) assign polarity (+/$-$). A key design choice is to force single-token outputs (e.g., ``C/N'', ``yes/no'') and use the model's log probabilities as ``model-implied'' confidence scores for each edge, enabling thresholding to trade off sparsity vs.\ coverage. To reduce prompt sensitivity, Theoraizer uses prompt ensembling (multiple prompt variants aggregated), and keeps humans in the loop by exposing intermediate tables/plots and encouraging iterative editing of variables and edges.

Methodologically, Theoraizer is closely aligned with CLD generation workflows that decompose structure construction into repeated pairwise decisions (e.g., prompting ``Does A cause B?''), and it explicitly targets scalability: evaluating $|V|(|V|-1)$ candidate edges becomes computational rather than purely cognitive labour. In validation, Theoraizer reports that LLM--human agreement on causal judgements is comparable to inter-expert agreement, suggesting LLMs can act as plausible ``stand-in'' raters for initial candidate CLDs, and supporting this by finding comparable LLM--human and human--human inter-rater agreement. At the same time, the paper emphasises limitations that mirror hallucination risks in causal structure generation: judgements are made pairwise and in isolation (encouraging dense graphs and ``short-circuit'' links that may actually be mediated), outputs reflect linguistic priors rather than epistemic truth, and biases/hallucinations necessitate expert review.

\subsubsection{Positioning of this thesis}

Contrasting their approach with this thesis, we add a mediating-variable check to address this limitation of reviewing links in isolation, and we apply hallucination detection and correction methods established in the broader field of LLM hallucination detection to CLD generation (e.g., LLM-as-a-judge, LLM-as-a-corrector, and uncertainty-quantification signals). In terms of using the logprobs-based uncertainty quantification, our approach slightly differs, as our generating model also produces a causal narrative through which variables are related (and which can also be used for post-hoc validation). Each token carries possible top-k alternatives over which the probabilities are distributed. This allows uncertainty quantification over this whole narrative, which is what this thesis will explore. All of this of course relies on the premise that these uncertainties reflect a model's tendency to hallucinate; this along with the hallucination reduction techniques used will be discussed later in this review.

\paragraph{Causal parrots: pattern retrieval versus genuine reasoning}
However, firstly we still want to address a broader question about what such edge-level pipelines recover from LLMs: are they primarily reproducing patterns that are well represented in training data, or do LLM outputs constitute causal reasoning? Work from the ``causal parrots'' hypothesis~\citep{zecevic2023causalparrots} offers a theoretical framework unifying these findings: LLMs encode causal facts within ``meta-SCMs'' (structural causal models describing relationships between causal facts in training data) but cannot perform genuine interventional or what-if reasoning. When LLMs succeed at causal tasks, they retrieve correlations between causal statements seen during training.

For our purpose, namely LLMs providing initial CLDs that reflect the literature, simply reproducing an accurate reflection of consensus on the causal relationships found in their literature may be enough. For genuine causal reasoning this model may break down and implies that models may accurately reproduce causal structures from well-documented domains that are part of training data while fabricating plausible-sounding but incorrect links in novel domains. The difficulty with this then becomes knowing what is part of the training data and what is not, which is often unclear, and also separate from the equally important questions of whether the model can faithfully reproduce that training data, or if the claims found in that data are true. All of these issues directly motivate hallucination mitigation approaches accompanying LLM-based CLD construction pipelines.

\section{How Hallucinations Manifest in Causal Structure Generation}

Hallucinations in CLD generation take distinct forms compared to free-text generation, and related work treats these failures through different benchmarks, annotations, and mitigation strategies. To synthesise this literature, we organise findings per hallucination type. Across the surveyed work, we summarise five recurring edge/graph-level error categories reported across benchmarks and expert analyses:

\textbf{Fabricated causal links} represent the most direct hallucination type---generating relationships between variables without textual grounding. The ReCAST benchmark~\citep{recast2025} directly evaluates multi-node causal graph construction from academic text---a task closely aligned with text-to-CLD causal translation. Unlike synthetic benchmarks, ReCAST uses unstructured, naturally written academic texts of varying length and complexity. Models must identify causal variables and construct directed graphs capturing the causal relationships described. The best-performing model achieves $F_1 = 0.477$, with common errors including fabricated links (e.g., ``loan usage flexibility $\rightarrow$ cash crop cultivation'' when no such relationship exists), difficulty with implicit causal information, and failure to connect causally relevant information spread across lengthy texts.

\textbf{Polarity errors} involve incorrect $+$/$-$ assignment to causal links. This manifests when LLMs state A increases B when it actually decreases B---a failure to apply ``ceteris paribus'' reasoning correctly. Unlike free-text hallucinations, CLDs require internally consistent sign patterns around feedback loops.

\textbf{Spurious variable introduction} creates variables not grounded in source text through over-generalisation. \citet{reinholtz2025pipelinealgebra} document examples like ``Macroeconomic and finance conditions'' with unknown polarity, or ``Global technology learning''---concepts too vague to support directional interpretation. Their study uses Pipeline Algebra to construct CLDs from text, with expert refinement identifying systematic LLM errors as hallucination mitigation.

\textbf{Missing causal links} represent faithfulness failures where explicit source relationships are omitted. The System Dynamics Bot study found that missing links often occur when ``authors assumed readers would know'' implicit information---highlighting the challenge of implicit versus explicit causal knowledge in text-to-CLD extraction.

\textbf{Direction reversal} and \textbf{DAG violations} constitute structural hallucinations unique to graph outputs. LLMs may reverse cause and effect or create cycles violating directed acyclic graph constraints. \citet{jiralerspong2024efficientcausalgraphdiscovery}, in their pairwise causal query approach to graph discovery, document that LLMs sometimes predict bidirectional causation or inconsistent edge directions when queried about the same variable pair in different orderings. Also in CLD construction, where feedback loops (or lack of them) have a crucial effect on the system's dynamics, this failure mode may be problematic.

\subsubsection{Formal CLD Hallucination Taxonomies}
\label{sec:sd_ai_taxonomy}

The \textbf{SD-AI benchmark}~\citep{schoenberg2025aibuildsdmodels} provides a systematic hallucination taxonomy for text-to-CLD tasks, maintaining a public leaderboard comparing model capabilities.\footnote{SD-AI Leaderboard: \url{https://ub-iad.github.io/sd-ai/\#/leaderboard/cld}} The taxonomy distinguishes: (1)~\textit{lack of conformance} (structural/formatting failures), (2)~\textit{variable extraction failures}, (3)~\textit{spurious edges} (false positives), (4)~\textit{missing edges} (false negatives), (5)~\textit{polarity errors}, and (6)~\textit{wrong causal narratives}. This thesis adopts categories 3--6 as the primary focus, with conformance addressed through structured outputs (JSON schema constraints) with the Judge being a qualitative check on conformance, and exploratory variable-level analysis relegated to Appendix~D, noting that main experiments use a ground truth given node set to isolate performance on edges to make the direct mapping to ground truths feasible by removing node-level differences from the comparison. The operational definitions and ground truth specifications used in our experiments appear in Section~\ref{sec:hallucination_definition}.

Expert evaluation at MIT~\citep{reinholtz2025pipelinealgebra} produced a complementary practical taxonomy with five categories: (A) ambiguous/overly broad variables, (B) duplicate/syntactically odd labels, (C) mis-specified linkage or polarity, (D) redundant elements, and (E) process-level limitations. These taxonomies provide actionable categories for automated detection systems in text-to-CLD pipelines, and although no comparable analysis for CLD generation approaches has been found (with closest work being the aforementioned), we adopt the SD-AI taxonomy to aid in making evaluation results directly comparable to text-to-CLD approaches.

\subsection{CLD Hallucination Mitigation Approaches}

Approaches that aim to reduce hallucinations through Retrieval-Augmented Generation show promise but remain constrained by retrieval quality. \citet{zhang2024causal} introduce LACR, an LLM-assisted causal discovery method that uses RAG over scientific papers to recover causal graphs: for each variable pair it retrieves relevant literature, prompts the LLM to infer (in)dependence and direct vs.\ indirect association in a constraint-based manner, aggregates these decisions to build the graph skeleton, and then uses the LLM to orient edges. Experiments on benchmark datasets show improved graph recovery and sensitivity to newer literature that may contradict outdated ``ground truth'' graphs.

Another approach integrates a causal graph into RAG. Causal-graph RAG (CausalRAG)~\citep{wang2025causalrag} shows promise for reducing hallucinations by retrieving causally relevant context beyond traditional RAG limitations like chunking, but they use the LLM for causal-path inference mainly post-retrieval, and either approach remains constrained by the coverage and quality of the underlying corpus/graph.

\subsubsection{RAG and the constraints of corpus dependence}
There are several failure modes when conditioning an LLM's generation on retrieved evidence. First, generation quality is fundamentally bounded by retrieval: the necessary evidence may be absent from the corpus (coverage/availability), inaccessible due to licensing or indexing constraints, or fragmented by preprocessing choices such as chunking and metadata loss. Second, even when the evidence exists, it may not be retrieved due to query--document mismatch (lexical or semantic mismatch, long-tail terminology, synonymy), embedding or retriever bias, suboptimal query formulation, or limitations of top-k recall under tight context budgets. Third, retrieval can return partially relevant or misleading context---e.g., semantically similar but causally irrelevant passages, outdated or contradictory sources, or redundant snippets that crowd out decisive evidence---so the model's response is conditioned on noise rather than signal. Finally, post-retrieval steps (reranking, summarization, graph/community aggregation) can further distort evidence by compressing away qualifiers and uncertainty, exacerbating omission and bias in the generated answer.

\subsubsection{Arguing for corpus independence from a practicality standpoint}
In CLD construction, if an expert must provide the corpus that is used as input to the CLD pipeline, the role of the model is constrained to a translation engine. The texts selected by the expert are presumably already, to some degree, familiar to the researcher, which can limit the model's ability to reduce bias---for example by suggesting hypotheses or search directions grounded in its broader training corpus that the researcher may be unfamiliar with. In contrast, if the model can be used to propose an initial causal hypothesis or search direction, this can save time and may yield new insight, conditional on the approach being reliable enough to improve discovery efficiency on average.

This motivates and positions corpus-independent hallucination reduction methods such as LLM-as-a-judge and logprobs uncertainty quantification, which are established in the broader hallucination reduction field but, according to the literature found in this review, remain untried in the context of causal loop diagram generation. Before discussing these methods in more detail in the next part of the literature review, we first conclude with an overview of the surveyed approaches.

\subsection{Comparison of LLM-Based CLD Approaches}

Table~\ref{tab:prior_cld_work} summarises prior work on LLM-based CLD construction, distinguishing approaches by their input requirements, hallucination mitigation strategies, and evaluation methods.

\begin{table}[H]
\centering
\caption[LLM-based CLD generation approaches]{Comparison of LLM-based CLD Generation and Validation Approaches}
\label{tab:prior_cld_work}
\footnotesize
\begin{tabular}{p{2.0cm}p{2.4cm}p{1.8cm}p{1.5cm}p{3.0cm}}
\toprule
\textbf{Approach} & \textbf{Method} & \textbf{Input} & \textbf{Corpus Req.} & \textbf{Hallucination Mitigation} \\
\midrule
SD-Bot~\citep{hosseinichimeh2024textmapdynamicsbot} & Text-to-CLD extraction via LLM & Text documents & Yes & None reported \\
\midrule
SD-AI~\citep{schoenberg2025aibuildsdmodels} & Benchmark suite (causal translation + conformance) & Text prompts / test cases & Yes & Evaluation benchmark (not a mitigation method) \\
\midrule
Theoraizer~\citep{Waaijers2024theoraizer} & Variable-pair querying; multiple candidates & User-defined variables & No & Logprob-based UQ; candidate generation \\
\midrule
Liu \& Keith~\citep{liu2025leveraging} & Curated prompting techniques & Dynamic hypothesis text & Yes & None reported \\
\midrule
Pipeline Algebra~\citep{reinholtz2025pipelinealgebra} & Typed operators; iterative prompting & Abstract + web search & Partial & Expert refinement step \\
\midrule
CausalRAG~\citep{wang2025causalrag} & RAG-enhanced causal discovery & Scientific literature & Yes & RAG grounding (99.5\% faithfulness) \\
\midrule
Susanti \& Holsm.~\citep{susanti2024prompting} & Validation (not generation); fine-tuning vs prompting & Existing causal graphs & Yes & Fine-tuning (+20.5 F1 over prompting) \\
\midrule
Zhang et al.~\citep{zhang2024causal} & RAG-based causal graph discovery & Scientific literature & Yes & RAG grounding \\
\midrule
\textbf{This thesis} & Corpus-independent generation + post-hoc validation & Domain spec.\ + GT Lit node set ($V$) & \textbf{No} & LLM-as-a-Judge; UQ metrics; Deep Research \\
\bottomrule
\end{tabular}
\end{table}

\noindent\textbf{Overall observations:}

\noindent Causal translation/extraction approaches outnumber corpus-independent generation-only approaches in the reviewed literature.

\begin{itemize}[noitemsep]
    \item \textbf{Corpus dependency:} Most approaches require an external corpus (text-to-CLD extraction or RAG), limiting applicability when relevant corpora are unavailable. Among the CLD-focused approaches surveyed in this review, Theoraizer~\citep{Waaijers2024theoraizer} and this thesis are notable for operating corpus-independently and for constructing CLD-level outputs, in contrast to adjacent pairwise causal discovery work that primarily targets causal graphs (often DAGs) rather than signed feedback-loop CLDs.
    \item \textbf{Hallucination mitigation:} Prior work largely neglects systematic hallucination detection outside of RAG, especially corpus-independent methods such as LLM-as-a-Judge. Theoraizer~\citep{Waaijers2024theoraizer} uses logprobs for uncertainty quantification but does not report a systematic evaluation of hallucination detection performance. RAG-based approaches (CausalRAG~\citep{wang2025causalrag}, Zhang et al.~\citep{zhang2024causal}) rely on retrieval grounding but are capped by retrieval quality. This thesis systematically evaluates multiple corpus-independent detection methods.
    \item \textbf{Evaluation scope:} SD-Bot~\citep{hosseinichimeh2024textmapdynamicsbot} reports 60\% link recall and 66--83\% loop recall on controlled datasets. Liu \& Keith~\citep{liu2025leveraging} demonstrate expert-comparable quality on simple structures. Pipeline Algebra~\citep{reinholtz2025pipelinealgebra} presents qualitative case studies. No prior work operationalises multiple definitions of LLM hallucinations in a CLD context (literature ground truth vs.\ synthetic CLD-specific hallucination types) and compares these systematically while controlling for prompt variation and LLM stochasticity.
\end{itemize}

\noindent\textbf{Positioning of our approach:} Given the prevalence of corpus-dependent text-to-CLD/causal translation pipelines and their reliance on retrieval quality, and the understudy of the (corpus-independent) alternative, CLD generation, combined with recent CLD-specific hallucination taxonomies, provides an opportunity for both quantifying hallucination rates of CLD generation and positioning the approach relative to causal translation. Also, this provides a platform to study whether established corpus-independent hallucination mitigation primitives---introduced later in this review, including LLM-as-a-Judge, LLM-as-a-Corrector, and logprob-derived uncertainty signals---not studied systematically in CLD generation, can reliably detect and remediate CLD hallucination types.

\noindent These CLD-specific failure modes and benchmark findings motivate an edge-level generation pipeline paired with post-hoc validation. The generator design evaluated in this thesis is described in Methods (Section~\ref{sec:generator}).

\noindent\textbf{Scope disclaimer.} While comprehensive ablation of the generator component is beyond the scope of this thesis, we report baseline performance metrics and share selected generator system experimentation conducted as part of the preliminaries to contextualize hallucination detection experiments; however, these results carry their own limitations as detailed in the appendix. It should also be noted that the generator design builds on intellectual contributions from Rick Quax, with implementation contributions from Cillian Hourican and Andrew Crossley, as well as contributions from Nitai Nijholt (as stated in the claim of authorship).
